\subsection{Циклы}

\textbf{Мотивирующий вопрос:} Нужно мигнуть светодиодом 100 раз. Будем копировать код 100 раз?

Циклы позволяют выполнять один и тот же блок кода многократно.

\subsubsection*{Цикл for}
Используется, когда известно количество повторений. Функция \texttt{range(N)} создает последовательность от 0 до N-1.

\subsubsection*{Цикл while}
Выполняется, пока условие истинно. Будьте осторожны: если условие никогда не станет ложным, получится <<бесконечный цикл>>.

\begin{codefile}[python]{examples/python/06_loops.py}
\end{codefile}

\begin{itemize}
    \item \texttt{break} — немедленно прерывает цикл.
    \item \texttt{continue} — переходит к следующей итерации.
\end{itemize}

\begin{taskbox}[Обратный отсчет]
Напишите программу, которая считает от 10 до 1, а затем пишет <<Поехали!>>. Используйте цикл \texttt{for} или \texttt{while}.
\end{taskbox}
